% ==================================================================
% id:       seq-009
% topic:    סדרות
% subtopic: אריתמטיקה של גבולות — הוכח או הפרך
% solution: full
% ==================================================================
\textbf{הוכיחו או הפריכו:}
\begin{enumerate}
    \item
    אם
    $(|a_n|)$
    מתכנסת, אז
    $(a_n)$
    חסומה.
\hint{היזכרו מה ידוע על חסימות של סדרה מתכנסת, ובדקו מה הקשר בין חסימות של $|a_n|$ לחסימות של $a_n$.}
\begin{solbox}
 נכון.
          למדנו כי אם סדרה מתכנסת אז היא חסומה, לכן מהנתון נסיק כי
          $|a_n|$
           היא חסומה, כלומר קיים
           $M>0$
           עבורו
           $|a_n|\le M$
           לכל
           $n$.
           מכאן נובע כי
           $-m\le a_n\le M$
           לכל
           $n$,
           וקיבלנו כי
           $a_n$
           חסומה.
\end{solbox}
    \item אם
    $(a_nb_n)$
    מתכנסת, אז שתי הסדרות
    $(a_n)$
    ו-
    $(b_n)$
    מתכנסות.
\hint{בדקו אם מכפלה יכולה להתכנס גם כששני הגורמים אינם מתכנסים. אם נראה לכם שהטענה לא נכונה, חפשו דוגמה נגדית פשוטה; אחרת נסו להוכיח אותה.}
\begin{solbox}
לא נכון.
נתבונן בסדרות
    $a_n=b_n=(-1)^n$.
    מתקיים
    $a_nb_n=1$
    לכל
    $n\in\mathbb{N}$,
    כלומר,
    הסדרה
    $a_nb_n$
    היא סדרה קבועה ולכן מתכנסת, אך הסדרות
    $a_n$
    ו-
    $b_n$
    אינן מתכנסות.
\end{solbox}
    %%%
    \item אם
    $(a_nb_n)$
    מתכנסת, אז לפחות אחת הסדרות
    $(a_n)$
    או
    $(b_n)$
    מתכנסת.
\begin{solbox}
        לא נכון. אותה דוגמה נגדית מסעיף קודם תעבוד גם כאן.
\end{solbox}
    \item
    אם
    $(a_n)$
    אינה חסומה אז בהכרח
    $a_n\to \infty$
    או
    $a_n\to-\infty$.
\hint{השוו בין מה שההגדרה של $a_n\to\infty$ דורשת לבין ההגדרה של סדרה לא חסומה. אם נראה לכם שהטענה לא נכונה, חפשו דוגמה נגדית פשוטה.}
\begin{solbox}
        לא נכון. הסדרה
        $a_n=n(-1)^n$
        אינה חסומה (לא מלעיל ולא מלרע) ומצד שני אינה מתכנסת לאינסוף או מינוס אינסוף.

$\bullet$
        דוגמה נגדית נוספת היא הסדרה
        $0,1,0,2,0,3,0,4,\ldots$
        שאינה חסומה מלעיל לכן אינה חסומה, וגם לא מתכנסת במובן הרחב כלל, בפרט לא ל-$\infty$
        ולא ל-$(-\infty)$.
\end{solbox}
    \item אם
    $(a_nb_n)$
    לא מתכנסת, אז
    $(a_n)$
    או
    $(b_n)$
    לא מתכנסות.
\hint{נסחו את הטענה השקולה לה (קונטרה-פוזיטיב) ובדקו אותה בעזרת אריתמטיקה של גבולות.}
\begin{solbox}
      נכון. איך נראה כי בהכרח אחת הסדרות אינה מתכנסת? נניח בשלילה כי שתיהן מתכנסות ונראה כי זה גורר סתירה.
      אכן, אם
      $a_n,b_n$
      שתיהן מתכנסות אז מאריתמטיקה נקבל כי גם
      $a_nb_n$
      מתכנסת, בסתירה מיידית לנתון.
\end{solbox}
    \item
    אם
    $a_n\to\infty$,
    אז
    $(a_n)$
    עולה.
\hint{השוו בין ההגדרה של $a_n\to\infty$ להגדרה של סדרה עולה. אם נראה לכם שהטענה לא נכונה, חפשו דוגמה נגדית פשוטה.}
\begin{solbox}
    לא נכון. נקח את הסדרה הכי פשוטה ששואפסת ל-$\infty$
    זו הסדרה
    $b_n=n$
    אך היא כמובן עולה, לכן ניצור ערבול קל בין איברי הסדרה שמצד אחד לא ישפיע על שאיפתה ל-$\infty$
    ומצד שני יהרוס את המונוטוניות:

    נתבונן בסדרה
    $.2,1,4,3,6,5,8,7,\ldots$
    זאת סדרה שגבולה הוא אינסוף
    אך היא אינה עולה.

$\bullet$
    אם חשוב לנו להציג סדרה כזו בעזרת נוסחה, להלן דוגמה נוספת:
    $$.a_n=n+(-1)^n$$
\end{solbox}
    \item
    אם
    $a_nb_n\to \infty$,
    אז
    $a_n\to\pm\infty$
    או
    $b_n\to\pm\infty$.
\hint{בדקו איך הגודל של $a_nb_n$ יכול להתחלק בין שתי הסדרות. אם נראה לכם שהטענה לא נכונה, חפשו דוגמה נגדית פשוטה.}
\begin{solbox}
לא נכון.
    בתור
    $(a_n)_{n=1}^\infty$
    נבחר את הסדרה הבאה:
    $$,10,1,1000,1,100000,1,\ldots$$
    ובתור
    $(b_n)_{n=1}^\infty$
    נבחר את הסדרה
    $$.1,100,1,10000,1,1000000,\ldots$$
    מתקיים
    $a_nb_n=10^n$,
    ולכן
    $a_nb_n\goes\infty$, אך אף אחת משתי הסדרות שהצגנו לא שואפת לא לאינסוף ולא למינוס אינסוף.

$\bullet$
דוגמה נגדית נוספת:
$$.a_n=b_n=(-1)^n\cdot n$$
\end{solbox}
    \item
    אם
    $(a_n),(b_n)$
    מתכנסות במובן הרחב, אז
    $(a_nb_n)$
    מתכנסת במובן הרחב.
\hint{עברו על המקרים האפשריים לגבולות של $a_n$ ו-$b_n$, ובדקו מה אומרת אריתמטיקה של גבולות במובן הרחב בכל אחד מהם.}
\begin{solbox}
        לא נכון.
        בהרבה מצבים זה נכון, למשל אם
        $a_n\goes\pm\infty$
        וגם
        $b_n\goes\pm\infty$
        או שאחת מהן מתכנסת למספר ששונה מ-$0$, למעט המצבים בהם סדרה אחת מתכנסת ל-$0$
        ואחת ל-$(\pm\infty)$, למשל
        $$a_n=n\goes\infty\qquad,\qquad b_n=\frac{(-1)^n}{n}\goes0$$
        כאן
        $a_nb_n=(-1)^n$
        אינה מתכנסת במובן הרחב.
\end{solbox}

    \item
    אם
    $(a_n)$
    חסומה ואם
    $b_n\to\infty$
    אז
    $a_n+b_n\to\infty$.
\hint{בדקו מה חסימות של $a_n$ נותנת לכם על $a_n+b_n$. אם נראה לכם שהטענה לא נכונה, חפשו דוגמה נגדית פשוטה.}
\begin{solbox}
        נכון. הסדרה
        $a_n$
        חסומה ולכן קיים
        $m\in\rr$
        עבורו
        $m\le a_n$
        לכל
        $n$.
        מכאן נובע כי
        $$m+b_n\le a_n+b_n$$
        ונזכיר כי מאריתמטיקה במובן הרחב אנו יודעים שמתקיים
        $\limn (m+b_n)=\infty$,
        לכן מכלל הפיצה גם
        $a_n+b_n\goes \infty$.
\end{solbox}
    \item
    אם
    $(a_n)$
    חסומה ואם
    $b_n\to\infty$
    אז
    $a_nb_n\to\infty$.

\hint{השוו לסעיף הקודם: האם חסימות של $a_n$ מספיקה גם כשמדובר במכפלה? אם נראה לכם שהטענה לא נכונה, חפשו דוגמה נגדית פשוטה.}
\begin{solbox}
        לא נכון. הסדרה הקבועה
        $a_n=0$
        היא חסומה והסדרה
        $b_n=n\goes\infty$
        מקיימות את הנתונים אך
        $a_nb_n=0$
        מתכנסת ל-$0$
        ולא ל-$\infty$.

        $\bullet$
        דוגמה נוספת היא
        $a_n=-1,b_n=n$
        כאן מתקיים
        $a_nb_n\goes -\infty$.

        $\bullet$
        דוגמה נוספת היא
        $a_n=(-1)^n,b_n=n$
        כאן מתקיים
        $a_nb_n=(-1)^nn$
        שאינה מתכנסת במובן הרחב כלל, בפרט לא ל-$\infty$.
\end{solbox}
\end{enumerate}
